\chapter{Common Instructions}

\section{Repository Structure}

Your Git repository must follow the naming convention below for
all module submissions. The evaluator's tooling (\texttt{pendragon})
relies on this structure to locate and test your work.

\begin{console}
<login>/
+-- <MODULE-ID>-<slug>/         # e.g. M-BLS-001-business-logic/
|   +-- <turnindir>/            # e.g. ex00-pricing-manipulation/
|   |   +-- <file1>             # e.g. poc-pricing.py
|   |   +-- <file2>             # e.g. report.md
|   |   +-- ...
|   +-- <turnindir>/            # e.g. ex01-race-conditions/
|   |   +-- ...
|   +-- ...
+-- ...
\end{console}

\begin{itemize}
  \item \textbf{Module directory}: the full directory name as shown in your
        curriculum (e.g. \texttt{M-BLS-001-business-logic}).
  \item \textbf{Turn-in directory}: the value of \texttt{\textbackslash turnindir}
        given in each exercise (e.g. \texttt{ex00-pricing-manipulation}).
  \item \textbf{Exercise files}: exactly the files listed in
        \texttt{\textbackslash exfiles}. No extra files, no missing files.
\end{itemize}

\section{What Constitutes a Valid Proof of Concept}

A PoC must be:

\begin{enumerate}
  \item \textbf{Reproducible}. Anyone with the same environment can run it
        and observe the vulnerability. Hard-coded paths, absolute IPs, and
        expired tokens are not acceptable.
  \item \textbf{Documented}. Every step is explained. The evaluator should
        not need to guess what the script does.
  \item \textbf{Targeted}. The exercise specifies an authorized target
        (lab, mock, local binary, or designated platform). Never test PoCs
        on production systems you do not own.
  \item \textbf{Minimal}. The PoC demonstrates \emph{only} the vulnerability
        in question. Unrelated payloads and noise are penalized.
\end{enumerate}

\begin{warning}
  A PoC that silently modifies system state, deletes files, or
  persists changes outside the exercise scope will be rejected.
  If your exploit requires cleanup, document the cleanup steps
  in your report.
\end{warning}

\section{Submission Workflow}

\begin{enumerate}
  \item Complete the exercise deliverables on your local machine.
  \item Commit and push to your assigned Git repository.
  \item The repository is frozen after the deadline.
  \item Peers clone your repository and evaluate using:
        \begin{itemize}
          \item \texttt{pendragon} for automated test execution
                (where the module supplies mouli configs).
          \item Manual rubric for report quality, code style, and
                methodology.
        \end{itemize}
\end{enumerate}

\hint{Always run \texttt{pendragon run <MODULE-ID> <ex\_id>
<path/to/student/repo>} before your final commit to verify
your submission passes the automated checks.}

\subsection*{Testing with pendragon}

If the module includes mouli configuration (JSON files in
\texttt{resources/mouli/}), you can test your submission locally:

\begin{console}
pendragon run M-BLS-001 ex00 ~/student-repos/tedsig42
\end{console}

\texttt{pendragon} resolves the exercise path automatically by
searching for \texttt{<mod-fullname>/<ex-dirname>/} under the
given student directory.

\subsection*{Style Checking with caliban}

For modules that require C code, \texttt{caliban} checks style
conformance:

\begin{console}
caliban check ~/student-repos/tedsig42/M-BLS-001-business-logic
\end{console}

\texttt{caliban} enforces the 42 Norm (function length, variable
naming, forbidden constructs) adapted for vulnerability research
contexts. Run it before submission --- failing style checks may
invalidate the exercise.

\section{No Extra Files Rule}

\begin{itemize}
  \item Your repository must contain \textbf{only} the files listed in
        \texttt{\textbackslash exfiles} for each exercise.
  \item Prohibited files: \texttt{.DS\_Store}, \texttt{\_\_pycache\_\_},
        \texttt{.venv/}, \texttt{node\_modules/}, \texttt{.git/},
        compiled binaries (unless explicitly listed in \texttt{exfiles}),
        editor swap files, log files, and any file not requested.
  \item Use \texttt{.gitignore} to exclude unwanted files automatically.
\end{itemize}

\section{Report Format}

When a report (\texttt{report.md}) is required, it must include:

\begin{enumerate}
  \item \textbf{Vulnerability class} (CWE identifier if applicable).
  \item \textbf{Reproduction steps} --- numbered, copy-paste ready.
  \item \textbf{Payloads} --- exact strings, requests, or inputs used.
  \item \textbf{Impact} --- what an attacker can achieve.
  \item \textbf{Mitigation} --- how the developer should fix it.
  \item \textbf{Sources} --- every resource consulted (see
        Section~\ref{sec:cite-sources}).
\end{enumerate}

\subsection*{Deliverable Header}

Every source file must include the deliverable header defined in
\path{../../../resources/latex/deliverable-header.en.tex}, adapted
to the language's comment syntax.

\section{Authorized Targets Only}

\begin{warning}
  You may \textbf{only} test against targets explicitly listed in
  the exercise (labs, local servers, mocks, designated binaries).
  Testing against real production systems, third-party applications,
  or any target you do not own or have written permission to test
  is prohibited and may result in immediate disqualification.
\end{warning}

Bug bounty ethics apply throughout this curriculum:
\begin{itemize}
  \item No testing on live programs without a valid scope agreement.
  \item No denial of service unless the exercise explicitly requires it.
  \item No data exfiltration beyond what is necessary to demonstrate
        the vulnerability.
  \item If you accidentally access real user data, stop immediately
        and document what happened.
\end{itemize}

\section{Cite All Sources}
\label{sec:cite-sources}

Any resource consulted during your work must be cited. This includes:

\begin{itemize}
  \item Academic papers and blog posts.
  \item Public bug bounty reports (HackerOne, Bugcrowd, etc.).
  \item Official documentation (RFCs, man pages, vendor docs).
  \item \textbf{AI tools}: model name, prompt summary, date, and
        verification method (see the AI Instructions chapter for
        the exact citation format).
  \item Code snippets adapted from open-source projects (with
        license attribution if required).
\end{itemize}

\subsection*{Citation format}

\begin{minted}{markdown}
# References

- [1] PortSwigger, "Business Logic Flaws",
  https://portswigger.net/web-security/logic-flaws
  (accessed 2025-03-15).
- [2] HackerOne Report #123456, "IDOR on /api/users",
  anonymized, accessed via public disclosure.
- [3] Generated with assistance from Claude 3.5 Sonnet
  on 2025-03-14. Prompt: "Write a Python script that...".
  Verified by: manual review of each HTTP request.
\end{minted}

\subsection*{Inline attribution}

If an exercise explicitly asks you to use a specific technique or
tool, you do not need to cite the exercise itself as a source.
Cite everything outside the exercise statement.
